4.0 PRZYGOTOWANIE DO GRY

Ogólna zasada:

Każdy z graczy powinien zaopatrzyć się w ołówek oraz jedną kartę cech dla każdej postaci, którą gra.

Żetony powinny zostać starannie porozdzielane oraz posegregowane według rodzajów.
Zaleca się, aby wszystkie zapisy na kartach cech dokonywane były ołówkiem, tak aby łatwo je było można zetrzeć, gdyż w trakcie gry cechy postaci ulegają ciągłym zmianom.

Zasady szczegółowe:

[4.1] Przygotowanie do gry musi przebiegać w następujący sposób:

A. WYBÓR BOHATERÓW.

1.
Postacie są podzielone na dwie grupy: bohaterowie oraz uczniowie.
Na żetonach, reprezentujących bohaterów znajduje się imię, zaś żetony, reprezentujące uczniów określają tylko ich rasę.

Żetony, reprezentujące bohaterów powinny zostać odwrócone zadrukowaną stroną ku dołowi oraz wymieszane tak, aby wybór był całkowicie przypadkowy.

2.
Jeżeli gra tylko jeden gracz wybiera on trzech bohaterów (przypadkowo) oraz trzech uczniów (dowolnie).

3.
Jeżeli w grze uczestniczy więcej niż jedna osoba kolejność w jakiej dokonywany jest wybór bohaterów ustalana jest przy pomocy rzutu kostką.
Wybór bohaterów dokonywany jest kolejno, zgodnie z ruchem wskazówek zegara, przy czym jako pierwszy dokonuje wyboru (,,ciągnie'') ten gracz, który wyrzucił największą liczbę oczek.

4.
Gracze ,,ciągną'' kolejno, aż do rozdzielenia między siebie trzech żetonów bohaterów i trzech żetonów uczniów.
Liczba postaci (tzn. bohaterów i uczniów) w grze może być mniejsza niż 6 (sugerowane minimum wynosi 4), jednak nie powinna przekraczać tej liczby.

B. OKREŚLENIE CECH POSTACI.

1.
Gracze odszukują cechy swoich bohaterów w tabeli 4.3 ,,Charakterystyka bohaterów'' i zapisują uzyskane w tej tabeli informacje na karcie cech danego bohatera.

2.
Gracze wybierają rasę swoich uczniów.
Uczeń może być krasnoludem, człowiekiem lub elfem.

3.
W zależności od rasy na karcie cech ucznia zapisywane są następujące informacje:

1 lub 2 oczka – dominuje czerwone słońce, liczba czarów określona jest lewą cyfrą współczynnika potencjał magiczny, 3 lub 4 oczka – dominuje żółte słońce – środkowa cyfra, 5 lub 6 oczek – dominuje niebieskie słońce – prawa cyfra.

Po ustaleniu dominującego słońca gracze wybierają z tych czarów, które zostały przypisane danej postaci w punkcie B6 taką ich liczbę, aby odpowiadała wielkości potencjału magicznego dla dominującego słońca.
Na przykład: jeżeli wyrzucone zostały 3 oczka, a potencjał magiczny postaci równy jest 4/5/6 gracz powinien wybrać (i zaznaczyć, np. przez podkreślenie) 5 spośród 6 posiadanych przez tę postać czarów.
Szósty, odrzucony czar podczas tej przygody nie będzie mógł być użyty.

E. SEGREGACJA ŻETONÓW.

Wszystkie żetony, reprezentujące komnaty i korytarze powinny zostać umieszczone w dużej filiżance lub innym tego typu pojemniku.
Żetony ,,Czarne Wrota'', ,,Brama labiryntu'' oraz ,,Bezimienni'' powinny zostać odłożone na bok.
Ich użycie jest określone specjalnymi przepisami.
Podobnie należy odłożyć na bok niewykorzystane żetony, reprezentujące bohaterów i uczniów.
W ciągu tej gry nie będą one już potrzebne.

F. ROZPOCZĘCIE GRY.

Żeton ,,Bramy labiryntu'' kładziony jest na stole, oznaczając tym samym miejsce, w którym znajduje się wejście.
Na tym żetonie powinien zostać położony żeton ,,Oddziału''.
W tym momencie gra może się rozpocząć.
Gracze przystępują do wykonywania czynności, podanych w części 5.0 ,,Przebieg gry''.

[4.2] Karta cech postaci – patrz plansza z tabelami. [4.3] Tabela ,,Charakterystyka bohaterów'' – patrz plansza z tabelami. [4.4] Tabela ,,Potencjał magiczny'' – patrz plansza z tabelami.